%% kapa: exam=1 round=37 year=2026 period=1 form=A subject=civil_law subject_ko=민법 seq=32 number=32
%% source: Q-Net 감정평가사 기출문제 (cst00309) artlSeq=5258924 file=raw/1008/5258924/2026년 제37회 감정평가사 1차 1교시.pdf
%% conversion: kordoc 4.24.51 -> tools/gen1.py (auto); review=proofread-sonnet; flags=none
\begin{question}{32}
관습상의 법정지상권에 관한 설명으로 옳은 것은? (다툼이 있으면 판례에 따름)
\begin{choices}
  \item 건물로서의 요건을 갖추었더라도 미등기의 무허가건물에 관해서는 관습상의 법정지상권이 인정될 수 없다.
  \item 관습상의 법정지상권이 성립한 후, 건물소유자가 토지소유자와 건물의 소유를 목적으로 하는 임대차계약을 체결하더라도 그 관습상의 법정지상권을 포기한 것으로 볼 수 없다.
  \item 타인의 토지 위에 그 토지소유자의 승낙을 얻어 신축한 건물을 매수한 경우, 매수인은 관습상의 법정지상권을 취득한다.
  \item 공유토지 위에 신축한 건물을 단독 소유하던 토지공유자 1인이 자신의 토지지분만을 양도하여 건물과 토지의 소유자가 달라진 경우에는 관습상의 법정지상권이 성립하지 않는다.
  \item 대지소유자가 그 지상건물을 다른 사람과 공유하면서 그 단독소유의 대지만을 제3자에게 매도한 경우, 건물공유자를 위한 관습상의 법정지상권은 인정되지 않는다.
\end{choices}
\end{question}
\answer{4}
